\section{V01 Introduction}

\begin{KR}{Agenda}
    \begin{itemize}
        \item Build Systems
        \begin{itemize}
            \item With/Without Build System
        \end{itemize}
        \item The Yocto Project
        \begin{itemize}
            \item With/Without Build System 
            \item Architecture
            \item Layers
            \item Recipes
            \item Adding an Application
            \item Configuring a System
            \item Building a bootable Image
        \end{itemize}
    \end{itemize}
\end{KR}

\begin{definition}{What is embedded Linux?}\\
    Embedded Linux is the usage of the Linux Kernel and various Open-Source Components in embedded systems.
\end{definition}

\begin{concept}{Advantages for Linux in Embedded}
\begin{itemize}
    \item Key advantage of Linux and open-source in embedded systems: Ability to re-use components
    \begin{itemize}
        \item Many components for network protocols, multimedia, graphic, cryptographic libraries, etc. exist
        \item Allows to quickly design and develop complicated products, based on existing components.
        \item No-one should re-develop yet another operating system kernel, TCP/IP stack, USB stack or another graphical toolkit library.
    \end{itemize}
    \item → Focus on the added value of your product.
    \item Low Cost! Free software can be duplicated on as many devices as you want, free of charge.
    \begin{itemize}
        \item → No Software License! (ercept for GPL) $\Rightarrow$ more details in lecture 09 
        \item Even the development tools are free, unless you choose a commercial embedded Linux edition.
    \end{itemize}
    \item Allows to have a higher budget for the hardware or to increase the company's skills and knowledge
\end{itemize}
\end{concept}

\begin{theorem}{Benefits of a build system}\\
\textbf{With a build system}
\begin{itemize}
    \item Build systems automate the process of building a target system, including the kernel.
    \item Automatic download, configure, compile and install all the components in the right order, sometimes after applying.
    \item They make sure all the application dependencies are matched.
    \item The build becomes reproducible, which allows to easily change the configuration of some components, upgrade them, fix bugs, etc.
\end{itemize}

\textbf{Without a build system}
    \begin{itemize}
        \item Each application has to be built manually
        \item The root file system has to be created from scratch.
        \item The applications configurations have to be done by hand.
        \item Each dependency has to be matched manually.
        \item Integrating software from different teams is painful 
    \end{itemize}
    \emoji{joy}
\end{theorem}

\begin{corollary}{Comparison of distribution projects}
    \begin{itemize}
        \item Buildroot
        \begin{itemize}
            \item Simple to use.
            \item Adapted for small embedded devices.
            \item Not perfect if you need advanced functionalities and multiple machines support.
            \item \href{http://buildroot.org/}{http://buildroot.org/}
        \end{itemize}
        \item OpenWRT
        \begin{itemize}
            \item Based on Buildroot.
            \item Primarily used for embedded network devices like routers.
            \item \href{http://openwrt.org/}{http://openwrt.org/}
        \end{itemize}
        \item Poky
        \begin{itemize}
            \item Part of the Yocto Project.
            \item Using OpenEmbedded.
            \item Suitable for more complex embedded systems.
            \item Allows lots of customization.
            \item Can be used for multiple targets at the same time.
            \item \href{http://yoctoproject.org/}{http://yoctoproject.org/}
        \end{itemize}
    \end{itemize}
\end{corollary}

\subsection{Yocto}

\begin{definition}{What is Yocto?}
    \begin{itemize}
        \item The Yocto Project is a set of templates, tools and methods that allow to build custom embedded Linux-based systems.
        \item It is an open source project initiated by the Linux Foundation in 2010 and is managed by one of its fellows: Richard Purdie.
    \end{itemize}
    
\end{definition}

